\subsection{The joined calculations retain separate validation targets}
\label{sec:joined-validation-targets}

The sector group in Section~\ref{sec:sector-company-composition} clarifies what
an end-to-end check must preserve. The bank's expected loan cash, its contractual
principal and coupon, the insurer's liability estimate, the parent's guarantee,
and each taxpayer's usable loss have different meanings. A group statement
can balance after one has been substituted for another. Its reference must
therefore retain the separate legal books, gross external cash and origin of
each adjustment, as well as the final totals.

Three kinds of comparison follow. For supplied economic and accounting facts,
an independently derived finite calculation tests the loan, tax, ownership,
currency and consolidation rules. Constant FX and unchanged ownership remove
one mechanism at a time: they test whether the calculation changes for the
right reason. For estimated loan or insurance cash, E4 instead tests the stated
conditional prediction against a credible tuned baseline. For numerical
valuation, E5 checks the declared measure and error in the downstream quantity.
A correct group bridge cannot establish the quality of a credit forecast or
insurance risk adjustment used as its input.

The same separation matters elsewhere. A DTA may be arithmetically correct for
supplied recoverable losses while its recovery forecast is unsupported. A POCI
schedule may conserve cash while using a forecast made after the reporting
date. A retrospective insurance method may follow its formula while using
hindsight that the route does not permit. Each comparison must therefore bind
the applicable law or policy, contractual and expected cash, historical
information vintage, valuation assumptions and ownership interval before
interpreting its result.

Some checks are exact within the stated domain: a hypothetical hedge option
must not generate actual cash; a component's notional funding must not become
new legal capital; and an NCI purchase after a close must not alter attribution
before that close. Other quantities, including predictive losses, tail estimates
and Monte Carlo accuracy, require their own uncertainty design. An error in the
target or reference invalidates the dependent comparison. A valid but poorly
performing forecast motivates the specified model repair. These outcomes
should lead to different next steps.

The original E1--E6 protocols consequently remain distinct after composition.
The newer examples add contract-loss order, dated credit receipts and write-offs,
dependent EPS, actual versus aligned option values, insurance transition
information and shared-sector attribution to their test populations. They
supply tractable checks before a more expensive run. Eligible data, paired
comparison budgets, dependence assumptions and decision-relevant margins
must still be fixed for each actual empirical target.
